% =============================================================================
%  RL-H2-01 | Finalisasi dan KUNCI protokol (v1)
% -----------------------------------------------------------------------------
%  Role    : RL — Research Lead
%  Sprint  : H2 | tenggat KERAS 17.00 | PARALEL
%  Ticket  : RL-H2-01                        Status: DONE | 2026-11-24 | protocol-v1
%  Target  : docs/eval_protocol_v1.tex  ->  docs/eval_protocol_v1.pdf
%  Input   : v0 + umpan balik MLE/CGS + definisi regime CGS-H2-04 + RL-H1-02
%  Output  : protokol terkunci + tag protocol-v1 + configs/global.yaml konsisten
%  Handoff : semua role (protokol = hukum sprint)
% =============================================================================
\documentclass[11pt,a4paper]{article}
% =============================================================================
%  00_preamble.tex | shared preamble untuk semua deliverable RL
%  Job     : bukan job card; dipakai oleh RL-H1-01 .. RL-H7-01
%  Target  : docs/ (salin bersama file .tex)
%  Catatan : kompilasi dengan pdflatex. Untuk XeLaTeX/LuaLaTeX, hapus inputenc.
% =============================================================================

\usepackage[utf8]{inputenc}
\usepackage[T1]{fontenc}
\usepackage{lmodern}
\usepackage[a4paper,margin=2.4cm]{geometry}
\usepackage{amsmath,amssymb}
\usepackage{booktabs}
\usepackage{longtable}
\usepackage{array}
\usepackage{graphicx}
\usepackage{xcolor}
\usepackage{enumitem}
\usepackage{float}
\usepackage{caption}
\usepackage{url}
\usepackage{microtype}
\usepackage[colorlinks=true,linkcolor=blue!50!black,urlcolor=blue!50!black,citecolor=blue!50!black]{hyperref}

% --- helper wajib -----------------------------------------------------------
% Path/nama file aman (underscore, garis miring) di dalam teks.
\newcommand{\fpath}[1]{\texttt{\detokenize{#1}}}
% TODO yang mencolok. Jangan dihapus sebelum job selesai.
\newcommand{\todo}[1]{\textcolor{red!75!black}{\textbf{[TODO: #1]}}}
\newcommand{\jobid}[1]{\textsf{\textbf{#1}}}
\newcommand{\runid}[1]{\texttt{\detokenize{#1}}}
\newcommand{\claim}[1]{\textsf{#1}}
\newcommand{\good}{\textcolor{green!45!black}{$\checkmark$}}
\newcommand{\bad}{\textcolor{red!75!black}{$\times$}}
\newcommand{\maybe}{\textcolor{orange!80!black}{$\sim$}}

\newcolumntype{L}[1]{>{\raggedright\arraybackslash}p{#1}}
\newcolumntype{C}[1]{>{\centering\arraybackslash}p{#1}}

\setlength{\parindent}{0pt}
\setlength{\parskip}{4pt}
\captionsetup{font=small,labelfont=bf}
\setlist[itemize]{leftmargin=1.2em,itemsep=2pt,topsep=2pt}
\setlist[enumerate]{leftmargin=1.5em,itemsep=2pt,topsep=2pt}

% Judul job yang konsisten di semua dokumen.
\newcommand{\jobheader}[4]{%
  \begin{center}
  {\Large\bfseries #1}\\[2pt]
  {\large #2}\\[4pt]
  \begin{tabular}{@{}ll@{}}
  \toprule
  Role & #3 \\
  Sprint & #4 \\
  \bottomrule
  \end{tabular}
  \end{center}
}


\title{\jobid{RL-H2-01} | Protokol Evaluasi v1.0 (TERKUNCI)\\
\large Sprint 1 --- Climate-Conditioned Dynamic Graph Learning}
\author{RENARO Lab \\ \small Research Lead (Alim)}
\date{24 November 2026 (H2) \quad | \quad \textcolor{green!45!black}{STATUS: TERKUNCI (\texttt{protocol-v1})}}

\begin{document}
\maketitle

\begin{center}
\fbox{\begin{minipage}{0.94\textwidth}
\textbf{Status: TERKUNCI} | Tanggal: 2026-11-24 (H2) | Disetujui: RL, MDE, MLE, CGS\\
Commit/tag: \texttt{protocol-v1} | File konfigurasi: \fpath{configs/global.yaml}\\
Perubahan setelah ini hanya lewat \fpath{docs/protocol_amendments.md}, dengan
alasan tertulis dan dicatat \emph{sebelum} hasil holdout diketahui.
\end{minipage}}
\end{center}

\section{Ringkasan perubahan dari v0 ke v1}
\begin{table}[H]
\centering
\small
\begin{tabular}{@{}L{3.8cm}L{4.8cm}L{5.8cm}@{}}
\toprule
Butir & v0 (Draf H1) & v1 (Final Terkunci H2) \\
\midrule
Horizon & 5, 10, 20; utama 10 & Dikunci: $h \in \{5, 10, 20\}$ hari kerja; utama $h = 10$ \\
Transformasi & $\log p$; metrik Rp/kg & Dikunci: target $\log P_{i,t+h}$; metrik di Rp/kg pasca $\exp(\hat{y})$ \\
Window input & 60 hari kerja & Dikunci: \texttt{window = 60} hari kerja (\texttt{panel\_start = 2017-07-01}) \\
CHIRPS availability & Opsi B utama, Opsi A sensitivitas & Dikunci dgn bukti HTTP HEAD probe (\jobid{CGS-H2-03}): Opsi B (\texttt{pentad\_end + 2d}) utama; Opsi A (\texttt{prelim\_sat + 5d} / \texttt{final\_sat} median 55{,}5d) sensitivitas \\
Regime: kering (\texttt{dry}) & $< P_{20}$, $\ge 3$ pentad & Dikunci (\jobid{CGS-H2-04}): \texttt{rain\_pentad\_mm < p20\_global} baseline pra-2018 selama $\ge 3$ pentad beruntun ($9.804$ node-days) \\
Regime: ekstrem (\texttt{extreme\_wet}) & $> P_{95}$ pentad & Dikunci (\jobid{CGS-H2-04}): \texttt{extreme\_5d == 1} ($> P_{95}$ baseline provinsi; rate $5{,}56\%$, $7.100$ node-days) \\
Prioritas label & $\text{dry} > \text{extreme} > \text{ENSO} > \text{normal}$ & Dikunci: $\texttt{dry} > \texttt{extreme\_wet} > \texttt{el\_nino} > \texttt{la\_nina} > \texttt{normal}$ (flag \texttt{CI lebar} jika $< 200$ node-days) \\
Ambang sintetis & Spearman $> 0{,}8$ / $> 0{,}7$ & Dikunci: Spearman $\rho > 0{,}80$ (1-edge), $\rho > 0{,}70$ (5-edge), MAE lag $\le 1{,}0$ pentad \\
Origin \& Holdout & O1--O3 + HOLDOUT & Dikunci: O1 (Test 2023), O2 (Test 2024), O3 (Test 2025--Mei 2026), HOLDOUT tersegel \texttt{2026-06-01..2026-10-31} (\jobid{CGS-H3-02}) \\
\bottomrule
\end{tabular}
\end{table}

\section{Nilai yang dibekukan (konsisten 100\% dengan \fpath{configs/global.yaml})}
\begin{table}[H]
\centering
\small
\begin{tabular}{@{}L{4.4cm}L{5.0cm}L{5.0cm}@{}}
\toprule
Field \fpath{configs/global.yaml} & Nilai Terkunci & Bukti / Dokumen \\
\midrule
\texttt{train\_start} & \texttt{2018-01-02} & Brief v2 \S7.2 (\texttt{panel\_start: 2017-07-01}) \\
\texttt{holdout\_start} & \texttt{2026-06-01} & Brief v2 \S7.2 \& \fpath{docs/holdout_seal.md} \\
\texttt{holdout\_end} & \texttt{2026-10-31} & \fpath{docs/holdout_seal.md} (\jobid{CGS-H3-02}) \\
\texttt{horizons} & \texttt{[5, 10, 20]} (\texttt{primary\_horizon: 10}) & \fpath{docs/decisions/ALL-H1-01.md} (\#1) \\
\texttt{embargo\_workdays} & \texttt{20} & \fpath{docs/decisions/ALL-H1-01.md} (\#4) \\
\texttt{seeds} & \texttt{[0, 1, 2, 3, 4]} & Brief v2 \S8 \\
\texttt{origins} & \texttt{O1, O2, O3, HOLDOUT} & \fpath{src/renaro/contract.py} (\texttt{ORIGINS}) \\
\texttt{window} & \texttt{60} & \S1 protokol ini \\
\texttt{regime\_definitions} & \texttt{p20\_global} ($\ge 3$ pentad), \texttt{p95} (\texttt{extreme\_5d}), prioritas 5-level, \texttt{min\_node\_days: 200} & \fpath{docs/regime_definitions.md} \& \fpath{reports/qa/regime_counts.md} (\jobid{CGS-H2-04}) \\
\texttt{chirps\_availability} & Opsi B utama (\texttt{pentad\_end + 2d}), Opsi A sensitivitas (\texttt{prelim + 5d}) & \fpath{reports/qa/availability.csv} \& \fpath{reports/qa/chirps_release_lag.md} (\jobid{CGS-H2-03}) \\
\bottomrule
\end{tabular}
\end{table}

\section{Definisi regime final (\jobid{CGS-H2-04})}
Label rezim di \fpath{data/interim/regime_labels.csv} berstatus \textbf{EVAL ONLY} dan dilarang keras dimasukkan ke \fpath{panel.parquet} atau matriks fitur model:
\begin{enumerate}
  \item \textbf{ENSO (\texttt{el\_nino}, \texttt{la\_nina}, \texttt{neutral})}: Berasal dari tabel episode RONI (\fpath{data/raw/enso/roni.csv}) menggunakan aturan NOAA 5-musim berturut-turut ($\ge +0{,}5^\circ\text{C}$ untuk El Ni\~no, $\le -0{,}5^\circ\text{C}$ untuk La Ni\~na). Setiap tanggal $t$ di-join secara \emph{as-of} ke musim terbaru dengan \texttt{available\_at <= t}.
  \item \textbf{Kekeringan (\texttt{dry})}: \texttt{rain\_pentad\_mm < p20\_global} selama $\ge 3$ pentad berturut-turut pada provinsi tersebut (\texttt{p20\_global} dihitung dari baseline pra-2018 provinsi tersebut).
  \item \textbf{Hujan Ekstrem (\texttt{extreme\_wet})}: \texttt{extreme\_5d == 1} (akumulasi pentad $> P_{95}$ baseline provinsi; frekuensi dasar $5{,}56\%$).
  \item \textbf{Normal (\texttt{normal})}: Tidak memenuhi \texttt{dry}, tidak memenuhi \texttt{extreme\_wet}, dan fase ENSO \texttt{neutral}.
  \item \textbf{Urutan Prioritas Label Utama}: $\texttt{dry} > \texttt{extreme\_wet} > \texttt{el\_nino} > \texttt{la\_nina} > \texttt{normal}$. Seluruh bendera biner individu tetap disimpan untuk evaluasi multi-label. Sel dengan $< 200$ \texttt{node\_days} diberi penanda \texttt{CI lebar}.
\end{enumerate}

\begin{table}[H]
\centering
\footnotesize
\begin{tabular}{@{}L{1.8cm}L{1.4cm}R{1.6cm}R{1.8cm}R{1.6cm}R{1.6cm}R{1.6cm}L{3.2cm}@{}}
\toprule
Origin & Part & \texttt{dry} & \texttt{extreme\_wet} & \texttt{el\_nino} & \texttt{la\_nina} & \texttt{normal} & Flag \texttt{CI lebar} ($<200$) \\
\midrule
O1 & train & 5.596 & 3.011 & 5.918 & 27.799 & 13.470 & --- \\
O1 & val   & 51    & 596   & 0     & 5.609  & 0      & \texttt{dry, el\_nino, normal} \\
O1 & test  & 1.990 & 402   & 2.341 & 5.669  & 2.008  & --- (semua $\ge 402$) \\
\midrule
O2 & train & 5.838 & 3.901 & 5.918 & 39.077 & 13.470 & --- \\
O2 & val   & 1.799 & 108   & 2.341 & 0      & 2.008  & \texttt{extreme\_wet, la\_nina} \\
O2 & test  & 1.066 & 797   & 4.463 & 1.880  & 4.238  & --- (semua $\ge 797$) \\
\midrule
O3 & train & 7.897 & 4.514 & 12.722 & 39.077 & 16.438 & --- \\
O3 & val   & 806   & 292   & 0      & 1.880  & 3.278  & \texttt{el\_nino} \\
O3 & test  & 661   & 1.585 & 0      & 12.518 & 2.780  & \texttt{el\_nino} \\
\midrule
HOLDOUT & train & 9.109 & 5.974 & 12.722 & 49.013 & 22.496 & --- \\
HOLDOUT & val   & 255   & 417   & 0      & 4.462  & 0      & \texttt{el\_nino, normal} \\
HOLDOUT & test  & \multicolumn{5}{c}{TERSEGEL (\texttt{2026-06-01..2026-10-31})} & Dibuka 1$\times$ di H6 \\
\midrule
\textbf{Total} & $< \text{Jun 26}$ & \textbf{9.804} & \textbf{7.100} & \textbf{12.722} & \textbf{58.311} & \textbf{28.955} & Total $116.892$ node-days \\
\bottomrule
\end{tabular}
\end{table}

\section{Aturan ketersediaan final (\jobid{CGS-H2-03})}
Tabel aturan ketersediaan rilis dari \fpath{reports/qa/availability.csv} (mencakup 24 fitur dalam 7 keluarga fitur di \fpath{docs/feature_dict.csv}):
\begin{table}[H]
\centering
\footnotesize
\begin{tabular}{@{}L{2.8cm}L{3.4cm}L{1.8cm}L{3.4cm}C{1.2cm}C{1.4cm}@{}}
\toprule
\texttt{feature\_family} & Sumber & Varian & Aturan (\texttt{rule}) & \texttt{lag\_days} & Pemutus \\
\midrule
\texttt{rain\_daily\_final} & CHIRPS v3.0 daily final/sat & \texttt{final\_sat} & \texttt{pentad\_end\_plus\_lag} & 2{,}0 & RL \\
\texttt{rain\_pentad\_final} & CHIRPS v3.0 daily final/sat & \texttt{final\_sat} & \texttt{pentad\_end\_plus\_lag} & 2{,}0 & RL \\
\texttt{rain\_prelim} & CHIRPS v3.0 daily prelim/sat & \texttt{prelim\_sat} & \texttt{release\_timestamp} & 5{,}0 & CGS \\
\texttt{roni} & CPC NOAA RONI (ERSSTv6) & \texttt{all} & \texttt{vintage\_or\_conservative\_3m} & --- & CGS \\
\texttt{production\_weight} & Satudata Pertanian / BPS & \texttt{all} & \texttt{year\_plus\_1} & --- & protocol \\
\texttt{calendar} & SKB 3 Menteri & \texttt{all} & \texttt{known\_in\_advance} & --- & protocol \\
\texttt{price} & PIHPS / SP2KP & \texttt{all} & \texttt{next\_workday} & 1{,}0 & protocol \\
\bottomrule
\end{tabular}
\end{table}

\noindent\textbf{Bukti probe HTTP HEAD (\fpath{reports/qa/chirps_release_lag_probe.csv}).}
Dari 6 berkas \texttt{final} (\texttt{2026-08-01..2026-10-31}), jeda rilis teramati adalah $18..109$ hari (median $55{,}5$ hari). Dari 6 berkas \texttt{prelim} (\texttt{2026-10-01..2026-11-20}), jeda rilis teramati adalah $2..8$ hari (median $4{,}0$ hari; 3 berkas terbaru tepat 2 hari). Hal ini memvalidasi keputusan \textbf{Opsi B (\texttt{pentad\_end + 2 hari})} sebagai proksi operasional utama dan \textbf{Opsi A (\texttt{prelim\_sat + 5 hari})} sebagai uji sensitivitas.

\section{Protokol evaluasi (mengikat)}
Seluruh bagian dari \fpath{docs/eval_protocol_v0.tex} berlaku penuh dan mengikat tanpa placeholder tersisa:
\begin{itemize}
  \item \textbf{Cakupan 34 provinsi (\jobid{MDE-H1-03})}: Berdasarkan \fpath{reports/qa/coverage_pihps.md}, 34 dari 34 provinsi memiliki cakupan $\ge 90\%$ di seluruh 8 tahun (2018--2025), sehingga seluruh 34 provinsi diikutkan.
  \item \textbf{Split \& Embargo (\S3 v0)}: Tiga rolling-origin (\texttt{O1}, \texttt{O2}, \texttt{O3}) dengan embargo 20 hari kerja di antara \texttt{train/val} dan \texttt{val/test}.
  \item \textbf{Metrik (\S6 v0)}: MAE (utama), RMSE, sMAPE dalam Rp/kg pada observasi asli (\texttt{is\_observed == True}), beserta \emph{Worst-Regime MAE} dan \emph{Regime Gap}.
  \item \textbf{Uji Statistik (\S7 v0)}: Uji Diebold--Mariano dua-sisi dengan varians Newey--West HAC ($L = h - 1$) dan koreksi Holm--Bonferroni ($\alpha = 0{,}05$) pada horizon utama $h = 10$.
  \item \textbf{RQ0 (\S8 v0)}, \textbf{Segel Holdout (\S9 v0, \texttt{2026-06-01..2026-10-31})}, \textbf{Verifikasi Sintetis (\S10 v0)}, \textbf{Perturbasi Prior 10--50\% (\S11 v0)}, dan \textbf{Pelaporan 5 Seed \texttt{0..4} (\S12 v0)} berlaku penuh.
\end{itemize}

\section{Amandemen}
\begin{table}[H]
\centering
\small
\begin{tabular}{@{}C{1.2cm}L{2.6cm}L{6.4cm}L{4.0cm}@{}}
\toprule
No & Tanggal & Alasan perubahan & Disetujui \\
\midrule
--- & --- & Belum ada amandemen pasca-penguncian \texttt{protocol-v1} (2026-11-24) & --- \\
\bottomrule
\end{tabular}
\end{table}

\section{Pernyataan kunci}
Dengan menandatangani dokumen ini, tim menyatakan:
\begin{enumerate}
  \item protokol dikunci sebelum model inti dijalankan pada data riil;
  \item tidak ada tuning pada periode uji;
  \item holdout hanya dibuka satu kali setelah Gate 2;
  \item hasil negatif akan dilaporkan.
\end{enumerate}

\vspace{6mm}
\begin{table}[H]
\centering
\small
\begin{tabular}{@{}L{2.4cm}L{4.4cm}L{3.8cm}L{3.0cm}@{}}
\toprule
Peran & Nama & Status Persetujuan & Tanggal \\
\midrule
RL  & Alim (\texttt{@alim}) & Disetujui \& Dikunci & 2026-11-24 \\
MDE & Akashuten (\texttt{@akashuten}) & Disetujui (H1 Data/Coverage) & 2026-11-24 \\
MLE & Agus Prasetya (\texttt{@agusprasetya}) & Disetujui (H1 Env/Loader) & 2026-11-24 \\
CGS & Muhammad Fariq Faqih (\texttt{@fariqfaqih}) & Disetujui (H1--H2 Climate/Seal) & 2026-11-24 \\
\bottomrule
\end{tabular}
\end{table}

\end{document}
